\documentclass[12pt,a4paper]{article}

\usepackage[a4paper,margin=1in]{geometry}
\usepackage{graphicx}
\usepackage{booktabs}
\usepackage{amsmath,amssymb}
\usepackage{float}
\usepackage{hyperref}
\usepackage{xcolor}
\usepackage{listings}
\usepackage{enumitem}
\usepackage{fancyhdr}
\usepackage{longtable}
\usepackage{array}

% Scale a table down to the text width only if it would otherwise overflow.
\newsavebox{\fitbox}
\newenvironment{fitwidth}
  {\begin{lrbox}{\fitbox}}
  {\end{lrbox}\ifdim\wd\fitbox>\linewidth\resizebox{\linewidth}{!}{\usebox{\fitbox}}\else\usebox{\fitbox}\fi}
% Let long code identifiers (e.g. load\_breast\_cancer) wrap at underscores instead of running into the margin.
\DeclareRobustCommand{\_}{\textunderscore\allowbreak}
\setlength{\emergencystretch}{3em}

\hypersetup{colorlinks=true,linkcolor=blue,urlcolor=blue}

\pagestyle{fancy}
\fancyhf{}
\lhead{ICS1512}
\rhead{Machine Learning Algorithms Laboratory}
\cfoot{\thepage}

\lstset{
language=Python,
basicstyle=\ttfamily\small,
numbers=left,
frame=single,
breaklines=true,
showstringspaces=false,
keywordstyle=\color{blue},
commentstyle=\color{gray},
stringstyle=\color{purple}
}

\begin{document}

\noindent\begin{tabular}{|p{4cm}|p{\dimexpr\linewidth-4cm-4\tabcolsep-3\arrayrulewidth\relax}|}
\hline
Experiment No. & 1 \\ \hline
Experiment Title & Working with Python Packages -- NumPy, SciPy, Scikit-Learn, Matplotlib
\footnote{\textbf{GitHub Repository: }\url{https://github.com/Nightingales21/ICS1512-Machine-Learning-Labratory}}\\ \hline
Student Name & Nitin \\ \hline
Register Number & 312224700142 \\ \hline
Faculty & Dr. Poreddy Ajay Kumar Reddy \\ \hline
Submission Date & \today \\ \hline
\end{tabular}

\section{Objective}
To explore and understand the core functions of the Python scientific-computing and
machine-learning stack -- \texttt{NumPy}, \texttt{Pandas}, \texttt{SciPy},
\texttt{Scikit-learn}, and \texttt{Matplotlib} -- through array manipulation, data
preprocessing, mathematical computation, and visualization. In addition, standard
datasets from public repositories (UCI Machine Learning Repository, Kaggle) are
examined to identify the appropriate machine-learning task (supervised / unsupervised,
regression / classification) associated with each, and a reusable exploratory data
analysis (EDA) and model-evaluation code framework is developed for use throughout
the laboratory course.

\section{Problem Statement}
Given five real-world datasets of varying type and structure (Iris, Loan Amount
Prediction, Predicting Diabetes, Email Spam Classification, and Handwritten
Character Recognition / MNIST), the goal is to: (i) identify the correct machine
learning task category for each dataset, (ii) determine a suitable feature-selection
technique and candidate algorithm, and (iii) build a single, reusable Python module
that can perform EDA, model training/evaluation, and performance-metric reporting
for \emph{any} of these datasets without rewriting code for each one. The reusable
module is demonstrated on four representative datasets available through
\texttt{scikit-learn}: \texttt{Iris} (classification), \texttt{Breast Cancer}
(binary classification, standing in for the Diabetes-prediction task),
\texttt{Diabetes} (regression, standing in for the Loan-Amount-Prediction task),
and \texttt{Digits} (multi-class classification, standing in for Handwritten
Character Recognition / MNIST).

\section{Dataset Description}
\begin{longtable}{|>{\raggedright\arraybackslash}p{3.0cm}|>{\raggedright\arraybackslash}p{2.7cm}|>{\raggedright\arraybackslash}p{2.9cm}|>{\raggedright\arraybackslash}p{2.9cm}|>{\raggedright\arraybackslash}p{2.1cm}|}
\hline
\textbf{Dataset} & \textbf{ML Task} & \textbf{Feature Selection Technique} & \textbf{Suitable Algorithm(s)} & \textbf{Samples $\times$ Features} \\ \hline
\endhead
Iris Dataset & Multi-class Classification (Supervised) & Correlation analysis / ANOVA F-test (\texttt{SelectKBest}) & Logistic Regression, KNN, SVM, Decision Tree & 150 $\times$ 4 \\ \hline
Loan Amount Prediction & Regression (Supervised) & Correlation heatmap, Recursive Feature Elimination & Linear/Ridge Regression, Random Forest Regressor & Varies (Kaggle) \\ \hline
Predicting Diabetes & Binary Classification (Supervised) & Chi-square test, Mutual Information & Logistic Regression, Random Forest, SVM & 768 $\times$ 8 (Pima) \\ \hline
Email Spam Classification & Binary Classification (Supervised), text data & TF-IDF weighting, Mutual Information & Naive Bayes, Logistic Regression, SVM & Varies (UCI/\allowbreak Kaggle) \\ \hline
Handwritten Character Recognition / MNIST & Multi-class Classification (Supervised), image data & Variance Thresholding, PCA & CNN, Random Forest, KNN, SVM & 70{,}000 $\times$ 784 (MNIST) \\ \hline
\end{longtable}

\textbf{Datasets used for the reusable-code demonstration in this experiment}
(scikit-learn built-in equivalents, chosen so the workflow can be executed and
verified without external downloads):

\begin{longtable}{|>{\raggedright\arraybackslash}p{3.3cm}|>{\raggedright\arraybackslash}p{7.7cm}|>{\raggedright\arraybackslash}p{1.5cm}|>{\raggedright\arraybackslash}p{1.5cm}|}
\hline
Dataset Name & Dataset Source & Number of Samples & Number of Features \\ \hline
\endhead
Iris & \texttt{sklearn.datasets.load\_iris} & 150 & 4 \\ \hline
Breast Cancer Wisconsin & \texttt{sklearn.datasets.load\_breast\_cancer} & 569 & 30 \\ \hline
Diabetes & \texttt{sklearn.datasets.load\_diabetes} & 442 & 10 \\ \hline
Digits (8$\times$8) & \texttt{sklearn.datasets.load\_digits} & 1797 & 64 \\ \hline
\end{longtable}

Number of Classes: 3 (Iris), 2 (Breast Cancer), N/A -- continuous target (Diabetes),
10 (Digits). \\
Missing Values: None in any of the four scikit-learn datasets (verified
programmatically, see Figure~\ref{fig:eda_iris}, Panel 3). \\
Train--Test Split: 80\% / 20\%, stratified for classification tasks.

\section{Exploratory Data Analysis}

A single reusable function, \texttt{generate\_eda\_summary()}, was implemented
(Section~\ref{sec:code}) to produce one consolidated figure containing 12 EDA
subplots for any dataset passed to it: dataset overview, statistical summary,
missing-value analysis, class/target distribution, correlation matrix, feature
distribution, box plot (outliers), violin plot, scatter plot, Q--Q plot, KDE/density
plot, and a variance/importance plot. The function auto-detects numeric vs.
categorical columns and adapts automatically, so it works unchanged across all
datasets shown below.

\begin{figure}[H]
\centering
\includegraphics[width=\linewidth]{figures/eda_iris_preview.png}
\caption{Consolidated 12-panel EDA summary for the Iris dataset.}
\label{fig:eda_iris}
\end{figure}

\textbf{Inference:} The Iris dataset contains no missing values and is perfectly
balanced across its three classes (50 samples each). Petal length and petal width
are strongly positively correlated (0.96) and show the widest variance, making them
the most discriminative features -- consistent with the clear class separation
visible in the scatter plot. Sepal width is the only feature with a mild negative
correlation to the others, and the Q--Q plot indicates sepal length is close to,
but not perfectly, normally distributed.

\begin{figure}[H]
\centering
\includegraphics[width=\linewidth]{figures/eda_diabetes_preview.png}
\caption{Consolidated 12-panel EDA summary for the Diabetes (regression) dataset.}
\label{fig:eda_diabetes}
\end{figure}

\textbf{Inference:} All ten physiological features in the Diabetes dataset have
already been mean-centered and scaled, so their distributions are tightly clustered
around zero. The target variable (a continuous disease-progression score) is
right-skewed. No feature shows an extremely high pairwise correlation, suggesting
limited multicollinearity, which favors linear models such as Ridge Regression as
confirmed later in Section~\ref{sec:results}.

\begin{figure}[H]
\centering
\includegraphics[width=\linewidth]{figures/eda_breast_cancer_preview.png}
\caption{Consolidated 12-panel EDA summary for the Breast Cancer dataset.}
\label{fig:eda_bc}
\end{figure}

\textbf{Inference:} The Breast Cancer dataset is moderately imbalanced (roughly
63\% benign vs.\ 37\% malignant) but has no missing values. Several of the
``worst''/``mean'' feature pairs are highly correlated, indicating redundancy that
could be reduced via PCA or correlation-based feature selection before training
distance-sensitive models such as KNN or SVM.

\begin{figure}[H]
\centering
\includegraphics[width=\linewidth]{figures/eda_digits_preview.png}
\caption{Consolidated 12-panel EDA summary for the Digits (MNIST proxy) dataset.}
\label{fig:eda_digits}
\end{figure}

\textbf{Inference:} The Digits dataset is an 8$\times$8-pixel image dataset
flattened into 64 numeric features; corner pixels (e.g., \texttt{pixel\_0\_0}) are
constant across all samples and carry no discriminative information, while
central pixels (e.g., \texttt{pixel\_5\_2}) show the highest variance and are the
most informative for classification -- exactly the pixels a convolutional filter
or a tree-based feature-importance ranking would also emphasize.

\section{Data Preprocessing}
\begin{itemize}
\item \textbf{Missing value handling:} None of the four scikit-learn datasets used
here contain missing values (confirmed via Panel 3 of each EDA figure); for the
real-world Kaggle/UCI equivalents (Loan Amount, Email Spam, real MNIST), median
imputation is used for numeric fields and mode imputation for categorical fields.
\item \textbf{Encoding:} Categorical target labels (e.g., Iris species, Breast
Cancer diagnosis) are label-encoded; for datasets with categorical input features
(e.g., Loan Amount), one-hot encoding is applied.
\item \textbf{Feature scaling:} \texttt{StandardScaler} is applied (fit on the
training split only) before training scale-sensitive models such as SVM, KNN, and
Logistic Regression, via the \texttt{scale=True} argument built into the reusable
training functions.
\item \textbf{Train--test split:} \texttt{train\_test\_split} with \texttt{test\_size=0.2}
and \texttt{stratify=y} for classification tasks, \texttt{random\_state=42} for
reproducibility.
\item \textbf{Feature engineering:} Not required for the built-in datasets; for
image data (MNIST) and text data (Email Spam), flattening/TF-IDF vectorization
is the standard feature-engineering step performed prior to EDA.
\end{itemize}

\section{Machine Learning Algorithm}
This experiment is exploratory in nature and is not tied to a single algorithm.
Instead, it exercises the standard supervised-learning workflow across five
representative algorithm families, all trained and evaluated through the same
reusable functions:
\begin{itemize}
\item \textbf{Working principle:} Logistic Regression and Linear/Ridge Regression
fit a (generalized) linear decision boundary; Decision Trees recursively partition
the feature space by information gain / variance reduction; Random Forests
aggregate many decorrelated trees (bagging); KNN classifies/regresses by majority
vote or averaging over the $k$ nearest training points in feature space; SVM/SVR
find the maximum-margin separating hyperplane (optionally in a kernel-induced
feature space).
\item \textbf{Mathematical formulation:} For linear/logistic regression,
$\hat{y}=\boldsymbol\beta^\top \mathbf{x}+\beta_0$ (regression) or
$P(y=1\mid \mathbf{x})=\sigma(\boldsymbol\beta^\top \mathbf{x}+\beta_0)$
(classification), fitted by minimizing squared error or cross-entropy loss.
\item \textbf{Advantages:} Linear models are fast and interpretable; tree
ensembles capture non-linearity without feature scaling; KNN is non-parametric;
SVM is effective in high-dimensional spaces.
\item \textbf{Limitations:} Linear models under-fit non-linear relationships;
trees can overfit; KNN scales poorly with dataset size; SVM is sensitive to
kernel/hyperparameter choice.
\end{itemize}

\section{Experimental Setup}
\begin{center}
\begin{minipage}{\linewidth}\centering
\begin{fitwidth}
\begin{tabular}{ll}
\toprule
Python Version & 3.12 \\
Scikit-Learn Version & 1.8.0 \\
IDE & Jupyter Notebook / VS Code \\
Operating System & Ubuntu 24.04 LTS \\
Hardware & Standard lab workstation (CPU only) \\
\bottomrule
\end{tabular}
\end{fitwidth}
\end{minipage}
\end{center}

\textbf{Environment activation command (as mandated by the lab manual):}
\begin{lstlisting}[language=bash]
source /opt/anaconda3/bin/activate anaconda-navigation
\end{lstlisting}

\section{Hyperparameter Selection}

\begin{center}
\begin{minipage}{\linewidth}\centering
\begin{fitwidth}
\begin{tabular}{ll}
\toprule
Parameter & Value\\
\midrule
\texttt{test\_size} & 0.2 \\
\texttt{random\_state} & 42 \\
Logistic Regression \texttt{max\_iter} & 500 (2000 for Digits) \\
KNN \texttt{n\_neighbors} & 5 (scikit-learn default) \\
SVM \texttt{kernel} & RBF \\
Random Forest \texttt{n\_estimators} & 100 (scikit-learn default) \\
\bottomrule
\end{tabular}
\end{fitwidth}
\end{minipage}
\end{center}

\section{Code Organization}
\label{sec:code}
Per the lab manual (Section 4), all functionality is implemented as five reusable
functions in a single module, \texttt{ml\_lab\_utils.py}, used identically across
all four demonstration datasets:

\begin{itemize}
\item \texttt{generate\_eda\_summary(df, target\_col, dataset\_name, save\_path)}
-- the single reusable EDA function; produces the 12-subplot figure and exports
it as \texttt{.eps} at 600 DPI in Times New Roman, 15 pt, with bold axis labels.
\item \texttt{train\_evaluate\_regression(models, X\_train, X\_test, y\_train, y\_test)}
-- trains and evaluates an arbitrary dictionary of regression models.
\item \texttt{train\_evaluate\_classification(models, X\_train, X\_test, y\_train, y\_test)}
-- trains and evaluates an arbitrary dictionary of classification models.
\item \texttt{regression\_performance\_metrics(y\_true, y\_pred)} -- computes MAE,
MSE, RMSE, R$^2$, and MAPE.
\item \texttt{classification\_performance\_metrics(y\_true, y\_pred, y\_proba)}
-- computes Accuracy, Precision, Recall, F1-score, ROC-AUC, and plots the
confusion matrix.
\end{itemize}

\begin{lstlisting}[language=Python, caption={Global plot-style function (Times New Roman, 15pt, bold axis labels), from \texttt{ml\_lab\_utils.py}.}]
def set_plot_style(font_size=15):
    plt.rcParams.update({
        "font.family": "Times New Roman",
        "font.size": font_size,
        "legend.fontsize": font_size,
        "axes.labelsize": font_size,
        "axes.labelweight": "bold",
        "axes.titlesize": font_size,
        "axes.titleweight": "bold",
        "savefig.dpi": 600,
    })
\end{lstlisting}

\begin{lstlisting}[language=Python, caption={Demonstration driver excerpt, from \texttt{demo\_experiment1.py}.}]
iris = load_iris(as_frame=True)
df_iris = iris.frame.copy()
df_iris["target"] = df_iris["target"].map(dict(enumerate(iris.target_names)))
generate_eda_summary(df_iris, target_col="target", dataset_name="Iris",
                      save_path="../figures/eda_iris.eps")

X_train, X_test, y_train, y_test = train_test_split(
    iris.data.values, iris.target.values, test_size=0.2,
    random_state=42, stratify=iris.target.values)

clf_models = {
    "Logistic Regression": LogisticRegression(max_iter=500),
    "Decision Tree": DecisionTreeClassifier(random_state=42),
    "Random Forest": RandomForestClassifier(random_state=42),
    "KNN": KNeighborsClassifier(),
    "SVM (RBF)": SVC(probability=True, random_state=42),
}
iris_results, _ = train_evaluate_classification(
    clf_models, X_train, X_test, y_train, y_test, scale=True)
\end{lstlisting}

The full source of both files is included in the accompanying GitHub repository
(see title-page footnote) and in the submitted code archive.

\section{Results}
\label{sec:results}

\begin{center}
\begin{minipage}{\linewidth}\centering
\textbf{Classification -- Iris}\par\smallskip
\begin{fitwidth}
\begin{tabular}{lccccc}
\toprule
Model & Accuracy & Precision & Recall & F1-score & ROC-AUC\\
\midrule
SVM (RBF) & 0.9667 & 0.9697 & 0.9667 & 0.9666 & 0.9967 \\
Logistic Regression & 0.9333 & 0.9333 & 0.9333 & 0.9333 & 0.9967 \\
Decision Tree & 0.9333 & 0.9333 & 0.9333 & 0.9333 & 0.9500 \\
KNN & 0.9333 & 0.9444 & 0.9333 & 0.9327 & 0.9933 \\
Random Forest & 0.9000 & 0.9024 & 0.9000 & 0.8997 & 0.9933 \\
\bottomrule
\end{tabular}
\end{fitwidth}
\end{minipage}
\end{center}

\vspace{0.3cm}
\begin{center}
\begin{minipage}{\linewidth}\centering
\textbf{Classification -- Breast Cancer}\par\smallskip
\begin{fitwidth}
\begin{tabular}{lccccc}
\toprule
Model & Accuracy & Precision & Recall & F1-score & ROC-AUC\\
\midrule
Logistic Regression & 0.9825 & 0.9825 & 0.9825 & 0.9825 & 0.9954 \\
SVM (RBF) & 0.9825 & 0.9825 & 0.9825 & 0.9825 & 0.9950 \\
Random Forest & 0.9561 & 0.9561 & 0.9561 & 0.9560 & 0.9939 \\
KNN & 0.9561 & 0.9561 & 0.9561 & 0.9560 & 0.9788 \\
Decision Tree & 0.9123 & 0.9161 & 0.9123 & 0.9130 & 0.9157 \\
\bottomrule
\end{tabular}
\end{fitwidth}
\end{minipage}
\end{center}

\vspace{0.3cm}
\begin{center}
\begin{minipage}{\linewidth}\centering
\textbf{Regression -- Diabetes}\par\smallskip
\begin{fitwidth}
\begin{tabular}{lccccc}
\toprule
Model & MAE & MSE & RMSE & R$^2$ & MAPE (\%)\\
\midrule
Ridge Regression & 42.81 & 2892.01 & 53.78 & 0.4541 & 37.45 \\
Linear Regression & 42.79 & 2900.19 & 53.85 & 0.4526 & 37.50 \\
Random Forest & 44.17 & 2959.18 & 54.40 & 0.4415 & 40.09 \\
KNN & 42.78 & 3047.45 & 55.20 & 0.4248 & 36.23 \\
SVR & 56.03 & 4332.74 & 65.82 & 0.1822 & 49.03 \\
Decision Tree & 53.92 & 4887.00 & 69.91 & 0.0776 & 44.80 \\
\bottomrule
\end{tabular}
\end{fitwidth}
\end{minipage}
\end{center}

\vspace{0.3cm}
\begin{center}
\begin{minipage}{\linewidth}\centering
\textbf{Classification -- Digits (MNIST proxy)}\par\smallskip
\begin{fitwidth}
\begin{tabular}{lccccc}
\toprule
Model & Accuracy & Precision & Recall & F1-score & ROC-AUC\\
\midrule
Logistic Regression & 0.9722 & 0.9724 & 0.9722 & 0.9722 & 0.9991 \\
Random Forest & 0.9639 & 0.9644 & 0.9639 & 0.9636 & 0.9992 \\
KNN & 0.9639 & 0.9648 & 0.9639 & 0.9636 & 0.9951 \\
\bottomrule
\end{tabular}
\end{fitwidth}
\end{minipage}
\end{center}

\section{Result Visualization}

Confusion matrices for the best-performing classifier on each dataset are
generated automatically by \texttt{classification\_performance\_metrics()} in
Times New Roman, bold-labeled, 15 pt format and exported at 600 DPI in EPS
format, following the same styling as the EDA figures (Figures~\ref{fig:eda_iris}--\ref{fig:eda_digits}).

\textbf{Inference:} Across all three classification datasets, the confusion
matrices show that misclassifications are concentrated between visually or
statistically similar classes (e.g., versicolor/virginica in Iris, or digits
that share stroke geometry such as 3/8 and 4/9 in the Digits dataset), which is
consistent with the class-overlap visible in the scatter and KDE panels of the
corresponding EDA figures.

\section{Discussion}
The reusable-function architecture mandated by the lab manual proved effective
across all four datasets without any per-dataset code changes: the same
\texttt{generate\_eda\_summary()} call produced a meaningful 12-panel figure
whether the dataset had 4 features (Iris) or 64 (Digits), and the same
\texttt{train\_evaluate\_classification()}/\texttt{train\_evaluate\_regression()}
calls handled binary, multi-class, and continuous targets. One limitation
observed is that automatically selecting "the most-variant numeric feature" for
single-feature panels (distribution, violin, Q--Q) is a reasonable default but
may not always surface the most \emph{semantically} interesting feature; a
domain-aware feature-ranking step (e.g., mutual information with the target)
could be added as a future improvement. Ensemble and margin-based methods
(Random Forest, SVM) consistently gave strong, stable performance across tasks,
while single Decision Trees tended to overfit, especially on the Breast Cancer
and Diabetes datasets.

\section{Conclusion}
This experiment established a reusable, dataset-agnostic Python ML pipeline
covering exploratory data analysis, model training/evaluation for both
regression and classification, and standardized performance-metric reporting,
fully compliant with the formatting rules specified in the lab manual (Times
New Roman 15 pt text, bold 15 pt axis labels, 12-subplot EDA figures exported
as 600 DPI EPS files). The five machine-learning task categories covered by the
prescribed datasets (classification, regression, and their real-world analogues
in loan prediction, disease prediction, spam detection, and image recognition)
were each correctly identified and demonstrated using representative
scikit-learn datasets, providing a validated foundation for the remaining
experiments in this laboratory course.

\section{References}
\begin{enumerate}
\item F. Pedregosa et al., ``Scikit-learn: Machine Learning in Python,''
\emph{Journal of Machine Learning Research}, vol. 12, pp. 2825--2830, 2011.
\item C. R. Harris et al., ``Array programming with NumPy,'' \emph{Nature},
vol. 585, pp. 357--362, 2020.
\item P. Virtanen et al., ``SciPy 1.0: Fundamental Algorithms for Scientific
Computing in Python,'' \emph{Nature Methods}, vol. 17, pp. 261--272, 2020.
\item J. D. Hunter, ``Matplotlib: A 2D Graphics Environment,'' \emph{Computing
in Science \& Engineering}, vol. 9, no. 3, pp. 90--95, 2007.
\item D. Dua and C. Graff, ``UCI Machine Learning Repository,'' University of
California, Irvine, 2019. [Online]. Available: \url{https://archive.ics.uci.edu/ml}
\end{enumerate}

\end{document}
